% SIAM Shared Information Template
% This is information that is shared between the main document and any
% supplement. If no supplement is required, then this information can
% be included directly in the main document.


% Packages and macros go here
\usepackage{lipsum}
\usepackage{amsfonts}
\usepackage{graphicx}
\usepackage{caption}
\usepackage{subfig}
\usepackage{epstopdf}
\usepackage{algorithmic}
\usepackage{booktabs}
\ifpdf
  \DeclareGraphicsExtensions{.eps,.pdf,.png,.jpg}
\else
  \DeclareGraphicsExtensions{.eps}
\fi

% Prevent itemized lists from running into the left margin inside theorems and proofs
\usepackage{enumitem}
\setlist[enumerate]{leftmargin=.5in}
\setlist[itemize]{leftmargin=.5in}

% Add a serial/Oxford comma by default.
\newcommand{\creflastconjunction}{, and~}

% Used for creating new theorem and remark environments
\newsiamremark{remark}{Remark}
\newsiamremark{hypothesis}{Hypothesis}
\newsiamremark{assumption}{Assumption}
\crefname{hypothesis}{Hypothesis}{Hypotheses}
\newsiamthm{claim}{Claim}
\newsiamremark{fact}{Fact}
\crefname{fact}{Fact}{Facts}

% Sets running headers as well as PDF title and authors
\headers{An Iterative Network for Parameter Estimation in Nonlinear Dynamical Systems from Sparse and Partial Observations}{Seong-heon Lee, Dongjin Lee, Jiaxi Gu, Joon Ha, and Jae-Hun Jung}

% Title. If the supplement option is on, then "Supplementary Material"
% is automatically inserted before the title.
\title{An Iterative Network for Parameter Estimation in Nonlinear Dynamical Systems from Sparse and Partial Observations\thanks{Submitted to the editors DATE.
\funding{This work was funded by the Fog Research Institute under contract no.~FRI-454.}}}

% Authors: full names plus addresses.
\author{
  Seong-Heon Lee\thanks{Graduate School of Artificial Intelligence, POSTECH, Pohang, Republic of Korea 
    (\email{shlee0125@postech.ac.kr}, \email{dongjinlee@postech.ac.kr}).}
  \and
  Dongjin Lee\footnotemark[2]
  \and
  Jiaxi Gu\thanks{Department of Mathematics, POSTECH, Pohang, Republic of Korea 
    (\email{jiaxigu@buffalo.edu}, \email{jung153@postech.ac.kr}).}
  \and
  Joon Ha\thanks{Department of Mathematics, Howard University, Washington, DC 20059, U.S.A. 
    (\email{joon.ha@Howard.edu}).}
  \and
  Jae-Hun Jung\footnotemark[3]
}

\usepackage{amsopn}
\DeclareMathOperator{\diag}{diag}

% ===== Hyperparameter highlighting (draft mode) =====
\newif\ifhpcolor
\hpcolortrue % draft: true, final submission: false

\ifhpcolor
  \usepackage{xcolor}
  \newcommand{\HP}[1]{\textcolor{blue}{#1}}
\else
  \newcommand{\HP}[1]{#1}
\fi

% ---------- Hyperparameter value macros ----------
\newcommand{\HPSys}{\HP{ogtt\_simul}}
\newcommand{\HPSeed}{\HP{42}}
\newcommand{\HPNumSamples}{\HP{10000}}
\newcommand{\HPAugFactor}{\HP{30}}
\newcommand{\HPTestSplit}{\HP{0.2}}
\newcommand{\HPBatchSize}{\HP{512}}

\newcommand{\HPEpochs}{\HP{10000}}
\newcommand{\HPLR}{\HP{1e-4}}
\newcommand{\HPWeightDecay}{\HP{1e-4}}
\newcommand{\HPEarlyStop}{\HP{True}}
\newcommand{\HPEarlyPat}{\HP{200}}
\newcommand{\HPEarlyDelta}{\HP{1e-6}}

\newcommand{\HPUseDeriv}{\HP{True}}
\newcommand{\HPDerivMethod}{\HP{spline}}
\newcommand{\HPIter}{\HP{10}}

\newcommand{\HPLossCfg}{\HP{supervised + recurrent (equal weights)}}

\newcommand{\HPHiddenDims}{\HP{[64,64,64]}}
\newcommand{\HPAct}{\HP{SiLU}}
\newcommand{\HPSN}{\HP{True}}

\newcommand{\HPSDEBias}{\HP{1.0}}
\newcommand{\HPSDEDiff}{\HP{1.27}}

%%% Local Variables: 
%%% mode:latex
%%% TeX-master: "ex_article"
%%% End: 
